\documentclass[10pt,a4paper]{amsart}
\usepackage{amsmath,amssymb,amsthm,graphicx}
% For commutative diagrams we use Xy-pic
\usepackage[all]{xy}

% We use 2cell for 2-commutative diagrams.
\xyoption{2cell}
% Theorem environments.
%
\theoremstyle{plain}
\newtheorem{theorem}[subsection]{Theorem}
\newtheorem{proposition}[subsection]{Proposition}
\newtheorem{lemma}[subsection]{Lemma}

\theoremstyle{definition}
\newtheorem{definition}[subsection]{Definition}
\newtheorem{example}[subsection]{Example}
\newtheorem{exercise}[subsection]{Exercise}
\newtheorem{situation}[subsection]{Situation}

\theoremstyle{remark}
\newtheorem{remark}[subsection]{Remark}
\newtheorem{remarks}[subsection]{Remarks}

\numberwithin{equation}{subsection}

% Macros
%
\def\lim{\mathop{\mathrm{lim}}\nolimits}
\def\colim{\mathop{\mathrm{colim}}\nolimits}
\def\Spec{\mathop{\mathrm{Spec}}}
\def\Hom{\mathop{\mathrm{Hom}}\nolimits}
\def\Ext{\mathop{\mathrm{Ext}}\nolimits}
\def\SheafHom{\mathop{\mathcal{H}\!\mathit{om}}\nolimits}
\def\SheafExt{\mathop{\mathcal{E}\!\mathit{xt}}\nolimits}
\def\Sch{\mathit{Sch}}
\def\Mor{\mathop{\mathrm{Mor}}\nolimits}
\def\Ob{\mathop{\mathrm{Ob}}\nolimits}
\def\Sh{\mathop{\mathit{Sh}}\nolimits}
\def\NL{\mathop{N\!L}\nolimits}
\def\CH{\mathop{\mathrm{CH}}\nolimits}
\def\proetale{{pro\text{-}\acute{e}tale}}
\def\etale{{\acute{e}tale}}
\def\QCoh{\mathit{QCoh}}
\def\Ker{\mathop{\mathrm{Ker}}}
\def\Im{\mathop{\mathrm{Im}}}
\def\Coker{\mathop{\mathrm{Coker}}}
\def\Coim{\mathop{\mathrm{Coim}}}

% Boxtimes
%
\DeclareMathSymbol{\boxtimes}{\mathbin}{AMSa}{"02}

%
% Macros for moduli stacks/spaces
%
\def\QCohstack{\mathcal{QC}\!\mathit{oh}}
\def\Cohstack{\mathcal{C}\!\mathit{oh}}
\def\Spacesstack{\mathcal{S}\!\mathit{paces}}
\def\Quotfunctor{\mathrm{Quot}}
\def\Hilbfunctor{\mathrm{Hilb}}
\def\Curvesstack{\mathcal{C}\!\mathit{urves}}
\def\Polarizedstack{\mathcal{P}\!\mathit{olarized}}
\def\Complexesstack{\mathcal{C}\!\mathit{omplexes}}
% \Pic is the operator that assigns to X its picard group, usage \Pic(X)
% \Picardstack_{X/B} denotes the Picard stack of X over B
% \Picardfunctor_{X/B} denotes the Picard functor of X over B
\def\Pic{\mathop{\mathrm{Pic}}\nolimits}
\def\Picardstack{\mathcal{P}\!\mathit{ic}}
\def\Picardfunctor{\mathrm{Pic}}
\def\Deformationcategory{\mathcal{D}\!\mathit{ef}}
\usepackage{cleveref}
\crefname{originaldefinition}{original definition}{original definitions}
\Crefname{originaldefinition}{Original definition}{Original definitions}
\makeatletter
\newlabel{situation-setup}{{02QL}{source}{}{sourcepointer-1}{}}
\newlabel{situation-setup@cref}{{[originaldefinition][0][]02QL}{[1][0][]source}{}{}{}}
\makeatother
\crefname{originalequation}{original equation}{original equations}
\Crefname{originalequation}{Original equation}{Original equations}
\makeatletter
\newlabel{homology-equation-cyclic-complex}{{02MW}{source}{}{sourcepointer-2}{}}
\newlabel{homology-equation-cyclic-complex@cref}{{[originalequation][0][]02MW}{[1][0][]source}{}{}{}}
\makeatother
\crefname{originallemma}{original lemma}{original lemmas}
\Crefname{originallemma}{Original lemma}{Original lemmas}
\makeatletter
\newlabel{lemma-length-finite}{{02QW}{source}{}{sourcepointer-3}{}}
\newlabel{lemma-length-finite@cref}{{[originallemma][0][]02QW}{[1][0][]source}{}{}{}}
\makeatother
\crefname{originaldefinition}{original definition}{original definitions}
\Crefname{originaldefinition}{Original definition}{Original definitions}
\makeatletter
\newlabel{definition-periodic-determinant}{{02PJ}{source}{}{sourcepointer-4}{}}
\newlabel{definition-periodic-determinant@cref}{{[originaldefinition][0][]02PJ}{[1][0][]source}{}{}{}}
\makeatother
\crefname{originalproposition}{original proposition}{original propositions}
\Crefname{originalproposition}{Original proposition}{Original propositions}
\makeatletter
\newlabel{proposition-length-determinant-periodic-complex}{{02QD}{source}{}{sourcepointer-5}{}}
\newlabel{proposition-length-determinant-periodic-complex@cref}{{[originalproposition][0][]02QD}{[1][0][]source}{}{}{}}
\makeatother
\title{For a cyclic two-periodic complex of a coherent sheaf with k-dimensional cohomology, the k-cycles of the two cohomology sheaves are rationally equivalent}
\author{The Stacks Project Authors}
\date{stacks-project:chow.tex:local-snapshot-20261012; GNU FDL 1.2 or later}
\begin{document}
\maketitle
\noindent\textit{Editorial note.} All delivered mathematical text is the author's own, copied from the pinned Stacks Project chapter file chow.tex; no statement or proof step was rewritten, reordered or removed.\par\medskip
\noindent\textit{Editorial note.} Preamble provenance: the pinned chapter file opens with an input of the shared preamble and carries no preamble text of its own, so the project's single Stacks-wide preamble staged at sources/stacks-chow/preamble.tex is excerpted verbatim as two spans, lines 14-18 (Xy-pic setup) and lines 154-217 (the native theorem declarations and mathematical macros). Nothing was added, removed or redefined; the wrapper class and the amsmath/amssymb/amsthm/graphicx base packages are supplied by the assembler.\par\medskip
\noindent\textit{Editorial note.} The counted statement and proof contain no cross reference, no citation and no included file, so the delivered body is self contained.\par\medskip
\noindent\textit{Editorial note.} Canonical location of the counted proposition: Stacks tag 02SB (sources/stacks-chow/pinned-original-source.tex line 18417).\par\medskip
\noindent\textit{Allowed standard prerequisites (editorial).}\begin{itemize}
\item The two-periodic complex attached to two endomorphisms of a module and its determinant. Homology, Equation labelled homology-equation-cyclic-complex (Stacks tag 02MW) and Same chapter, Definition labelled definition-periodic-determinant (source line 16611), invoked at statement line 18431 and proof lines 18456-18461. Fix the complex notation (F, phi, psi), its cohomology sheaves, and the determinant of the complex at the generic point of a component.
\item Length of modules with finite dimensional support and local finiteness of the top dimensional components. Same chapter, Lemma labelled lemma-length-finite (source line 1702), invoked at proof line 18452. Supplies that the family of components of dimension k+1 is locally finite, so that the sum of divisors is a well defined cycle.
\item The length-determinant relation for a two-periodic complex over a one dimensional local ring. Same chapter, Proposition labelled proposition-length-determinant-periodic-complex (source line 18024), invoked at proof line 18505. Supplies that the e-invariant of the localised complex equals the sum of the orders of the localised determinants, which is the coefficient comparison closing the proof; it is retained only as a tagged navigation pointer.
\item The base situation (S, delta) and the rational equivalence of k-cycles. Same chapter, Situation labelled situation-setup (source line 1372), invoked in the statement. Fixes delta-dimensions, the attached k-cycles and the relation \textasciitilde{} rat.
\end{itemize}\medskip
\ifdefined\hypertarget\hypertarget{sourcepointer-1}{}\fi
\noindent\textit{Source reference.} Original definition 02QL: sources/stacks-chow/pinned-original-source.tex, source lines 1372--1372. Same-chapter pointer into the author environment situation at source line 1372; invoked at proof line None of the counted proof. It is retained only as a navigation pointer; every load-bearing step is inside the delivered span.\par\medskip
\ifdefined\hypertarget\hypertarget{sourcepointer-2}{}\fi
\noindent\textit{Source reference.} Original equation 02MW: sources/stacks-chow/tags.txt, source lines 3459--3459. Cross-chapter standard prerequisite located in the project tag list staged in the local snapshot; invoked at proof line None of the counted proof. It is retained only as a navigation pointer; every load-bearing step is inside the delivered span.\par\medskip
\ifdefined\hypertarget\hypertarget{sourcepointer-3}{}\fi
\noindent\textit{Source reference.} Original lemma 02QW: sources/stacks-chow/pinned-original-source.tex, source lines 1702--1702. Same-chapter pointer into the author environment lemma at source line 1702; invoked at proof line 18452 of the counted proof. It is retained only as a navigation pointer; every load-bearing step is inside the delivered span.\par\medskip
\ifdefined\hypertarget\hypertarget{sourcepointer-4}{}\fi
\noindent\textit{Source reference.} Original definition 02PJ: sources/stacks-chow/pinned-original-source.tex, source lines 16611--16611. Same-chapter pointer into the author environment definition at source line 16611; invoked at proof line 18461 of the counted proof. It is retained only as a navigation pointer; every load-bearing step is inside the delivered span.\par\medskip
\ifdefined\hypertarget\hypertarget{sourcepointer-5}{}\fi
\noindent\textit{Source reference.} Original proposition 02QD: sources/stacks-chow/pinned-original-source.tex, source lines 18024--18024. Same-chapter pointer into the author environment proposition at source line 18024; invoked at proof line 18505 of the counted proof. It is retained only as a navigation pointer; every load-bearing step is inside the delivered span.\par\medskip
\begin{lemma}
\label{lemma-maps-between-coherent-sheaves}
Let $(S, \delta)$ be as in Situation \ref{situation-setup}.
Let $X$ be a scheme locally of finite type over $S$.
Let $\mathcal{F}$ be a coherent sheaf on $X$.
Let
$$
\xymatrix{
\ldots \ar[r] &
\mathcal{F} \ar[r]^\varphi &
\mathcal{F} \ar[r]^\psi &
\mathcal{F} \ar[r]^\varphi &
\mathcal{F} \ar[r] & \ldots
}
$$
be a complex as in Homology, Equation (\ref{homology-equation-cyclic-complex}).
Assume that
\begin{enumerate}
\item $\dim_\delta(\text{Supp}(\mathcal{F})) \leq k + 1$.
\item $\dim_\delta(\text{Supp}(H^i(\mathcal{F}, \varphi, \psi))) \leq k$
for $i = 0, 1$.
\end{enumerate}
Then we have
$$
[H^0(\mathcal{F}, \varphi, \psi)]_k
\sim_{rat}
[H^1(\mathcal{F}, \varphi, \psi)]_k
$$
as $k$-cycles on $X$.
\end{lemma}

\begin{proof}
Let $\{W_j\}_{j \in J}$ be the collection of irreducible
components of $\text{Supp}(\mathcal{F})$
which have $\delta$-dimension $k + 1$. Note that $\{W_j\}$
is a locally finite collection of closed subsets of
$X$ by Lemma \ref{lemma-length-finite}.
For every $j$, let $\xi_j \in W_j$ be the generic point.
Set
$$
f_j = \det\nolimits_{\kappa(\xi_j)}
(\mathcal{F}_{\xi_j}, \varphi_{\xi_j}, \psi_{\xi_j})
\in
R(W_j)^*.
$$
See Definition \ref{definition-periodic-determinant} for notation.
We claim that
$$
- [H^0(\mathcal{F}, \varphi, \psi)]_k + [H^1(\mathcal{F}, \varphi, \psi)]_k
=
\sum (W_j \to X)_*\text{div}(f_j)
$$
If we prove this then the lemma follows.

\medskip\noindent
Let $Z \subset X$ be an integral closed subscheme of $\delta$-dimension $k$.
To prove the equality above it suffices to show that the coefficient $n$
of $[Z]$ in
$
[H^0(\mathcal{F}, \varphi, \psi)]_k - [H^1(\mathcal{F}, \varphi, \psi)]_k
$
is the same as the coefficient $m$ of $[Z]$ in
$
\sum (W_j \to X)_*\text{div}(f_j)
$.
Let $\xi \in Z$ be the generic point.
Consider the local ring $A = \mathcal{O}_{X, \xi}$.
Let $M = \mathcal{F}_\xi$ as an $A$-module.
Denote $\varphi, \psi : M \to M$ the action of $\varphi, \psi$ on
the stalk.
By our choice of $\xi \in Z$ we have $\delta(\xi) = k$
and hence $\dim(\text{Supp}(M)) = 1$.
Finally, the integral closed subschemes
$W_j$ passing through $\xi$ correspond to the minimal primes
$\mathfrak q_i$ of $\text{Supp}(M)$.
In each case the element $f_j \in R(W_j)^*$ corresponds to
the element $\det_{\kappa(\mathfrak q_i)}(M_{\mathfrak q_i}, \varphi, \psi)$
in $\kappa(\mathfrak q_i)^*$. Hence we see that
$$
n = - e_A(M, \varphi, \psi)
$$
and
$$
m =
\sum
\text{ord}_{A/\mathfrak q_i}
(\det\nolimits_{\kappa(\mathfrak q_i)}(M_{\mathfrak q_i}, \varphi, \psi))
$$
Thus the result follows from
Proposition \ref{proposition-length-determinant-periodic-complex}.
\end{proof}

\end{document}
